%% ============================================================
%%  MorphNN — Research Paper (Methodology, Results, Discussion)
%%  Target: Q1 Elsevier / Nature-level journal
%%  Sections: Methodology · Results · Discussion
%%  Author will add: Title, Abstract, Keywords, Introduction,
%%                   Related Work, Conclusion, References
%% ============================================================
\documentclass[preprint,12pt]{elsarticle}

\usepackage{lineno,hyperref}
\modulolinenumbers[5]

%% Mathematics
\usepackage{amsmath,amssymb,amsfonts,bm}

%% Algorithms
\usepackage{algorithm}
\usepackage{algpseudocode}
\usepackage{algorithmicx}
\algrenewcommand\algorithmicindent{1em}

%% Tables
\usepackage{booktabs}
\usepackage{multirow}
\usepackage{multicol}
\usepackage{tabularx}
\usepackage{array}
\usepackage{makecell}
\usepackage{rotating}
\usepackage{siunitx}
\sisetup{detect-all}

%% Figures
\usepackage{graphicx}
\usepackage{subcaption}
\usepackage{xcolor}

%% Misc
\usepackage{url}
\usepackage[hidelinks]{hyperref}

\journal{Computers in Biology and Medicine}   % change as needed

\begin{document}

\begin{frontmatter}

\title{MorphNN: A Multimodal Morphological-Visual Framework for
       Viral Exanthem Triage with Clinically Interpretable Predictions}

%% Authors — fill in
\author[1]{[Author~1]}
\author[1]{[Author~2]}
\address[1]{[Department, University, City, Country]}
\cortext[cor1]{Corresponding author. E-mail: [email]}

\begin{abstract}
[To be added by authors.]
\end{abstract}

\begin{keyword}
multimodal medical AI \sep viral exanthem classification \sep
morphological feature engineering \sep joint deep learning fusion \sep
clinical error reduction \sep gradient-weighted class activation mapping
\end{keyword}

\end{frontmatter}

\linenumbers

%% ===========================================================
%%  NOTE: Sections 1 (Introduction) and 2 (Related Work)
%%        will be added by the authors.
%%        Numbering begins at 3 below; adjust if needed.
%% ===========================================================


%% ===========================================================
\section{Methodology}
\label{sec:methodology}
%% ===========================================================

We propose \textbf{MorphNN}, a multimodal deep learning framework that
jointly fuses two complementary diagnostic modalities: (i)~a visual
modality comprising raw dermoscopic or clinical images, processed by a
pre-trained convolutional backbone, and (ii)~a structured morphological
modality comprising 16 deterministic clinical descriptors extracted from
the same image.  The framework is trained end-to-end, enabling both
modalities to co-adapt through a shared classification objective.  An
overview is illustrated in Fig.~\ref{fig:pipeline}.

To isolate the contribution of each design decision, we evaluate
MorphNN within a four-phase ablation that progressively integrates
visual and morphological evidence.  The ablation phases are: (1)~\textbf{Morph-RF}
--- morphological features alone with a Random Forest classifier;
(2)~\textbf{FrozenCNN+M} --- frozen CNN embeddings concatenated with morphological
features, classified by RF; (3)~\textbf{CNN-FT} --- a fine-tuned CNN
without morphological features (the strongest purely visual baseline);
(4)~\textbf{HybridRF} --- the original MorphNN hybrid strategy, combining
PCA-compressed frozen CNN features with morphological features in an RF
classifier.  \textbf{MorphNN} (the proposed architecture) supersedes all
ablation phases via joint end-to-end multimodal training.


% -----------------------------------------------------------
\subsection{Dataset Description}
\label{sec:datasets}
% -----------------------------------------------------------

\subsubsection{Multi-Class Viral Skin Lesion Dataset (MCVSLD) ---
               Primary Evaluation}
\label{sec:mcvsld}

The primary dataset is the Multi-Class Viral Skin Lesion Dataset
(MCVSLD)~\cite{mcvsld2023}, comprising clinical and dermatoscopic images
of six viral exanthem conditions: Chickenpox, Cowpox,
Hand-Foot-and-Mouth Disease (HFMD), Healthy skin, Measles, and
Monkeypox.  The dataset contains 11,325 images sourced from clinical
archives and public repositories, partitioned into 9,060 training,
1,131 validation, and 1,134 held-out test images
(Table~\ref{tab:dataset}).  The validation partition was not used
separately; cross-validation was applied internally to the training
partition, as detailed in Section~\ref{sec:setup}.

The training partition exhibits significant class imbalance, with
Monkeypox representing the majority class (3,408 images, $37.6\%$)
and Measles the minority (660 images, $7.3\%$).  To reduce this
imbalance without modifying the test distribution, geometric
augmentations (random horizontal flip, rotation within $\pm 20^\circ$,
and colour jitter: $\Delta$brightness$= 0.2$, $\Delta$contrast$= 0.2$,
$\Delta$saturation$= 0.1$, $\Delta$hue$= 0.05$) were applied
exclusively within the training partition to bring underrepresented
classes to a minimum of 2,000 samples, yielding 13,408 augmented
training images.  The test set was never augmented, preserving the
authentic class frequency for evaluation.

\subsubsection{HAM10000 Skin Lesion Dataset ---
               Cross-Domain Validation}
\label{sec:ham}

To assess cross-domain generalisability, we evaluated MorphNN on the
Human Against Machine with 10,000 training images
(HAM10000)~\cite{tschandl2018ham10000}, a benchmark of 10,015
dermoscopic images across seven categories: actinic keratoses
(\textit{akiec}), basal cell carcinoma (\textit{bcc}), benign
keratosis-like lesions (\textit{bkl}), dermatofibroma (\textit{df}),
melanoma (\textit{mel}), melanocytic nevi (\textit{nv}), and vascular
lesions (\textit{vasc}).  The dataset is severely imbalanced: \textit{nv}
constitutes $66.9\%$ of all images ($n = 6{,}705$), while \textit{df}
and \textit{vasc} together contribute fewer than 260 images.

A distinguishing feature of our evaluation protocol is
\emph{lesion-level} partitioning.  The HAM10000 dataset includes
multiple photographs of the same lesion, identified via a
\texttt{lesion\_id} field.  Many published evaluations perform
image-level splits, allowing photographs of the same lesion to appear
in both train and test sets, which constitutes evaluation leakage.
We partitioned all unique lesions stratified by diagnosis in an
$80/20$ ratio, yielding 8,020 training images and 1,995 test images
that are guaranteed free of within-lesion overlap.  The training set
was augmented identically to MCVSLD, reaching 17,368 samples.
Per-class counts are given in Table~\ref{tab:dataset}.

\textit{Important note on morphological feature alignment.}
The 16 MorphNN descriptors were specifically designed for viral
exanthems, which present as \emph{distributed multi-lesion rashes}
(Chickenpox: many discrete vesicles; Monkeypox: fewer, larger
pustules; HFMD: acral distribution).  HAM10000 images are
\emph{dermoscopic close-ups of a single focal lesion}.  This
domain mismatch affects five of the 16 features
(Section~\ref{sec:feature_alignment}); the remaining 11 features
(shape regularity, colour heterogeneity, lesion dimensions) retain
meaningful signal for single-lesion classification.  Cross-domain
results are thus reported as a generalisability stress test, not a
competitive benchmark against lesion-specific models.

\begin{table}[htbp]
\centering
\caption{Per-class image counts for both datasets.
  Training counts are reported before and after augmentation.
  Test counts reflect the unaugmented natural distribution.}
\label{tab:dataset}
\setlength{\tabcolsep}{5pt}
\begin{tabular}{lrrr|rr}
\toprule
\multirow{2}{*}{\textbf{Class}} &
  \multicolumn{3}{c|}{\textbf{MCVSLD (viral exanthem)}} &
  \multicolumn{2}{c}{\textbf{HAM10000 (dermoscopy)}} \\
\cmidrule(lr){2-4}\cmidrule(lr){5-6}
 & Train (orig.) & Train (aug.) & Test &
   Train (orig.) & Test \\
\midrule
Chickenpox / \textit{akiec} & 900    & 2{,}000 & 113 & 262  & 62   \\
Cowpox / \textit{bcc}       & 792    & 2{,}000 & 99  & 411  & 100  \\
HFMD / \textit{bkl}         & 1{,}932 & 2{,}000 & 242 & 882  & 218  \\
Healthy / \textit{df}       & 1{,}368 & 2{,}000 & 171 & 92   & 28   \\
Measles / \textit{mel}      & 660    & 2{,}000 & 83  & 891  & 218  \\
Monkeypox / \textit{nv}     & 3{,}408 & 3{,}408 & 426 & 5{,}364 & 1{,}337 \\
{---} / \textit{vasc}       & ---    & ---     & --- & 114  & 32   \\
\midrule
\textbf{Total}              & 9{,}060 & 13{,}408 & 1{,}134 & 8{,}016 & 1{,}995\\
\bottomrule
\end{tabular}
\end{table}


% -----------------------------------------------------------
\subsection{Algorithms in Detail}
\label{sec:algorithms}
% -----------------------------------------------------------

\subsubsection{Morphological Feature Extraction: A Deterministic
               16-Descriptor Pipeline}
\label{sec:morph_features}

Morphological feature extraction operates on raw BGR images prior to any
learned transformation.  The pipeline is fully deterministic,
parameter-free, and requires no training data, making it interpretable
by design.  Given an input image
$\mathbf{I} \in \mathbb{R}^{H \times W \times 3}$, extraction proceeds
in four stages.

\paragraph{Stage 1: Colour-Space Conversion and Lesion Segmentation.}
The image is converted to the HSV colour space,
$\mathbf{I}_{HSV} = \phi_{BGR\to HSV}(\mathbf{I})$.  Lesions in viral
exanthems exhibit elevated saturation relative to surrounding healthy
skin.  An adaptive binary mask is obtained via Otsu's
method~\cite{otsu1979} on the saturation channel $S$:
\begin{equation}
  M(u,v) =
  \begin{cases}
    1 & \text{if } S(u,v) \geq \tau^*,\\
    0 & \text{otherwise,}
  \end{cases}
  \quad
  \tau^* = \underset{\tau}{\arg\max}\;\sigma_B^2(\tau),
  \label{eq:otsu}
\end{equation}
where $\sigma_B^2(\tau)$ is the between-class variance at threshold
$\tau$.  The mask is refined by morphological closing (disk radius 5)
to fill intraregional holes and opening (disk radius 2) to remove
noise pixels.

\paragraph{Stage 2: Connected Component Analysis.}
Lesion instances are extracted via 8-connectivity analysis on $M$:
\begin{equation}
  \mathcal{L} = \{L_1, L_2, \ldots, L_N\},\quad N = |\mathcal{L}|,
  \label{eq:components}
\end{equation}
where $L_i$ is the pixel set of the $i$-th connected region.  Regions
with fewer than 50 pixels are discarded.  For each $L_i$, the following
properties are computed: area $A_i = |L_i|$, perimeter $P_i$, bounding
box $(w_i, h_i)$, and convex hull area $\hat{A}_i$.

\paragraph{Stage 3: Colour Statistics on Lesion and Background.}
Let $\mathcal{P}_L = \bigcup_i L_i$ denote lesion pixels and
$\mathcal{P}_B = \{(u,v):M(u,v)=0\}$ background pixels.  Per-channel
vectors $\mathbf{h}_L$, $\mathbf{s}_L$, $\mathbf{v}_L$ (lesion HSV) and
$\mathbf{s}_B$ (background saturation) are assembled for feature
computation.

\paragraph{Stage 4: Feature Vector Assembly.}
Table~\ref{tab:morph_features} defines all 16 descriptors constituting
$\mathbf{f} \in \mathbb{R}^{16}$.  Features $f_1$--$f_6$ and $f_{14}$,
$f_{15}$ quantify lesion burden and distribution (multi-lesion rash
properties); $f_4$, $f_7$--$f_{13}$, $f_{16}$ quantify shape and colour
properties that are informative for both multi-lesion rashes and
single-lesion dermoscopy.  After extraction, $\mathbf{f}$ is standardised
using a StandardScaler fitted on the training split.

\begin{table}[htbp]
\centering
\caption{The 16 deterministic morphological descriptors of MorphNN.
  Features marked $\dagger$ are designed for multi-lesion distributed
  rashes and have limited transferability to single-focal-lesion
  dermoscopic datasets (e.g.\ HAM10000).}
\label{tab:morph_features}
\renewcommand{\arraystretch}{1.15}
\begin{tabular}{clp{6.5cm}}
\toprule
$f_k$ & \textbf{Name} & \textbf{Definition} \\
\midrule
$f_1$$^\dagger$ & Lesion count &
  $N = |\mathcal{L}|$ \\
$f_2$ & Average lesion area &
  $\frac{1}{N}\sum_i A_i$ \\
$f_3$ & Area heterogeneity &
  $\sqrt{\frac{1}{N}\sum_i(A_i - f_2)^2}$ \\
$f_4$ & Average circularity &
  $\frac{1}{N}\sum_i \frac{4\pi A_i}{P_i^2}$\quad($=1$: perfect circle) \\
$f_5$$^\dagger$ & Sparsity score &
  $1 - \frac{\sum_i A_i}{H \cdot W}$ \\
$f_6$$^\dagger$ & Confluence density &
  $\frac{|\phi_\delta(\mathcal{P}_L)| - \sum_i A_i}{\sum_i A_i}$,\;
  $\phi_\delta$ = dilation (radius 10) \\
$f_7$ & Localised hue &
  $\overline{\mathbf{h}_L}$ \\
$f_8$ & Localised saturation &
  $\overline{\mathbf{s}_L}$ \\
$f_9$ & Average aspect ratio &
  $\frac{1}{N}\sum_i \frac{w_i}{h_i}$ \\
$f_{10}$ & Average solidity &
  $\frac{1}{N}\sum_i \frac{A_i}{\hat{A}_i}$ \\
$f_{11}$ & Localised value &
  $\overline{\mathbf{v}_L}$ \\
$f_{12}$ & Hue spread &
  $\operatorname{std}(\mathbf{h}_L)$ \\
$f_{13}$ & Saturation spread &
  $\operatorname{std}(\mathbf{s}_L)$ \\
$f_{14}$$^\dagger$ & Spatial entropy &
  $-\sum_{x,y} p(x,y)\log_2 p(x,y)$,\; $p$ = 2-D KDE of lesion
  centroids \\
$f_{15}$$^\dagger$ & Max-area ratio &
  $\frac{\max_i A_i}{\sum_i A_i}$ \\
$f_{16}$ & Background saturation &
  $\overline{\mathbf{s}_B}$ \\
\bottomrule
\multicolumn{3}{l}{\footnotesize $\dagger$ Limited transferability to
single-lesion dermoscopic images; see Section~\ref{sec:feature_alignment}.}
\end{tabular}
\end{table}

These features capture clinical morphological properties familiar to
dermatologists: lesion burden ($f_1$--$f_3$), shape regularity ($f_4$,
$f_9$, $f_{10}$), spatial distribution ($f_5$, $f_6$, $f_{14}$,
$f_{15}$), and colour characteristics ($f_7$, $f_8$, $f_{11}$--$f_{13}$,
$f_{16}$).  For example, Chickenpox presents with a high lesion count
($f_1$) and near-circular vesicles ($f_4 \approx 1$), Monkeypox with
fewer, larger, and more irregular pustules ($f_2 \uparrow$,
$f_3 \uparrow$), and HFMD with a characteristic acral distribution
reflected in spatial entropy ($f_{14}$).


\subsubsection{Morph-RF: Morphology-Only Baseline (Phase~1)}
\label{sec:phase1}

Morph-RF establishes a baseline using only $\mathbf{f}$ as input to a
class-balanced Random Forest:
\begin{equation}
  \hat{y} = \mathrm{RF}(\mathbf{f};\;\theta),\quad
  \theta = \{\text{200 trees, balanced class weights,
  max-features} = \sqrt{16}\}.
  \label{eq:phase1}
\end{equation}
Class-balanced weights $w_c = N/(K \cdot N_c)$ are used throughout all
RF-based configurations.  This phase quantifies the discriminative power
of morphological descriptors in isolation, without any pixel-level
visual information.


\subsubsection{FrozenCNN+M: Frozen Transfer Learning with
               Morphological Augmentation (Phase~2)}
\label{sec:phase2}

FrozenCNN+M combines frozen deep visual embeddings with the
morphological vector.  Seven CNN backbone architectures
$\mathcal{B} = \{\text{MobileNetV2, MobileNetV3-L, ShuffleNetV2,
SqueezeNet1.1, ResNet-18, ResNet-34, EfficientNet-B0}\}$ are evaluated.
For each backbone $b$, a visual embedding is extracted via global
average pooling (GAP) over the last feature map:
\begin{equation}
  \mathbf{x}_{vis} = \mathrm{GAP}\!\left(g_b(\mathbf{I})\right) \in \mathbb{R}^{d_b},
  \label{eq:phase2_emb}
\end{equation}
with embedding dimension $d_b \in \{1280, 960, 1024, 512, 512, 512, 1280\}$
for MobileNetV2 through EfficientNet-B0, respectively.
Backbone weights remain fixed at ImageNet-pretrained values.
The fused vector $[\mathbf{x}_{vis};\,\mathbf{f}] \in \mathbb{R}^{d_b + 16}$
is classified by a class-balanced RF.


\subsubsection{CNN-FT: Fine-Tuned CNN Standalone Baseline (Phase~3)}
\label{sec:phase3}

CNN-FT is the primary competitive reference: each backbone $b \in \mathcal{B}$
is fine-tuned end-to-end on the visual training data, replacing the
ImageNet head with a $K$-class linear classifier.  Training employs
AdamW~\cite{loshchilov2019adamw} with cosine annealing warm
restarts~\cite{loshchilov2017sgdr}:
\begin{equation}
  \eta(t) = \eta_{\min} +
  \tfrac{1}{2}(\eta_{\max} - \eta_{\min})
  \left(1 + \cos\!\left(\tfrac{t\bmod T_0}{T_0}\pi\right)\right),
  \label{eq:cosine_annealing}
\end{equation}
with $\eta_{\max} = 10^{-4}$, $\eta_{\min} = 10^{-6}$, $T_0 = 5$.
The loss is label-smoothing weighted cross-entropy~\cite{muller2019labelsmoothing}:
\begin{equation}
  \mathcal{L}_{LS} = -\sum_{c=1}^K w_c\,\tilde{y}_c \log\hat{p}_c,\quad
  \tilde{y}_c = (1-\varepsilon)y_c + \tfrac{\varepsilon}{K},\quad
  \varepsilon = 0.05.
  \label{eq:label_smoothing}
\end{equation}
Morphological features are not used in CNN-FT; it represents the
strongest possible single-modality visual baseline.


\subsubsection{HybridRF: Original Morphological-Visual Hybrid (Phase~4)}
\label{sec:phase4}

HybridRF compresses frozen CNN embeddings via PCA to six components,
then concatenates with the morphological vector for RF classification:
\begin{equation}
  \mathbf{x}_{\mathrm{fused}} =
  \left[\mathbf{f};\;\mathrm{PCA}_6\!\left(\mathbf{x}_{vis}\right)\right]
  \in \mathbb{R}^{22}.
  \label{eq:phase4_fusion}
\end{equation}
This phase tests whether na\"ive post-hoc concatenation is sufficient
to capture multimodal complementarity, without any joint gradient flow.


\subsubsection{MorphNN: Proposed Joint Multimodal Architecture}
\label{sec:morphnn}

\textbf{MorphNN} is the proposed architecture, designed to overcome the
fundamental limitation of all preceding phases: the absence of
cross-modal gradient flow.  In MorphNN, both modalities are trained
jointly through a single differentiable objective, enabling the visual
branch to adapt its representations in light of morphological evidence
and vice versa (Fig.~\ref{fig:architecture}).

\paragraph{Visual branch.}
The convolutional stack of MobileNetV3-Large~\cite{howard2019mobilenetv3}
($g : \mathbb{R}^{3\times 224 \times 224} \to \mathbb{R}^{960\times7\times7}$,
initialised with ImageNet weights) extracts a 960-dimensional visual embedding:
\begin{equation}
  \mathbf{x}_{vis} = \mathrm{GAP}\!\left(g(\mathbf{I})\right) \in \mathbb{R}^{960}.
  \label{eq:morphnn_visual}
\end{equation}

\paragraph{Morphological branch.}
The standardised morphological vector $\hat{\mathbf{f}} \in \mathbb{R}^{16}$
is projected to a 64-dimensional embedding:
\begin{equation}
  \mathbf{x}_{morph} =
  \mathrm{Dropout}_{0.3}\!\left(\mathrm{Hardswish}\!\left(
    \mathbf{W}_m\hat{\mathbf{f}} + \mathbf{b}_m\right)\right)
  \in \mathbb{R}^{64},\quad
  \mathbf{W}_m \in \mathbb{R}^{64\times16}.
  \label{eq:morphnn_morph}
\end{equation}
Hardswish activation ($\mathrm{hs}(x) = x\,\mathrm{ReLU6}(x+3)/6$)
is used for consistency with the backbone.

\paragraph{Multimodal fusion and classification head.}
The embeddings from both branches are concatenated to form a joint
multimodal representation:
\begin{equation}
  \mathbf{z} = [\mathbf{x}_{vis};\;\mathbf{x}_{morph}] \in \mathbb{R}^{1024},
  \label{eq:morphnn_fusion}
\end{equation}
which is passed through a two-layer classification head:
\begin{equation}
  \hat{\mathbf{p}} = \mathrm{softmax}\!\left(
    \mathbf{W}_2\cdot\mathrm{Dropout}_{0.2}\!\left(
      \mathrm{Hardswish}\!\left(\mathbf{W}_1\mathbf{z} + \mathbf{b}_1\right)
    \right) + \mathbf{b}_2
  \right),
  \label{eq:morphnn_head}
\end{equation}
where $\mathbf{W}_1 \in \mathbb{R}^{512\times1024}$,
$\mathbf{W}_2 \in \mathbb{R}^{K\times512}$.

\paragraph{Differential learning rates and training protocol.}
The key training design choice is modality-specific learning rates:
\begin{equation}
  \eta_{\text{visual backbone}} = 10^{-4},\qquad
  \eta_{\text{morph branch + head}} = 10^{-3}.
  \label{eq:lr_differential}
\end{equation}
The lower rate for the visual backbone preserves the structured
representations acquired during ImageNet pretraining, while the higher
rate for the morphological branch and fusion head allows rapid
adaptation to disease-specific patterns.  The shared AdamW optimiser
(weight decay $\lambda = 10^{-4}$) with cosine annealing warm restarts
($T_0 = 10$, Eq.~\eqref{eq:cosine_annealing}) is applied for 20 epochs,
batch size 32, using the label-smoothing class-weighted loss
(Eq.~\eqref{eq:label_smoothing}).  Algorithm~\ref{alg:morphnn_train}
provides the complete training procedure.

\begin{algorithm}[htbp]
\caption{MorphNN Joint Multimodal Training}
\label{alg:morphnn_train}
\begin{algorithmic}[1]
\Require Images $\{\mathbf{I}_n\}_{n=1}^{N_{tr}}$,
         morph vectors $\{\mathbf{f}_n\}$, labels $\{y_n\}$,
         epochs $E{=}20$, batch size $B{=}32$
\State Initialise visual backbone $g$ with ImageNet weights
\State Initialise $\mathbf{W}_m$, $\mathbf{W}_1$, $\mathbf{W}_2$ with Kaiming uniform
\State Fit StandardScaler $\rightarrow$ $\hat{\mathbf{f}}_n = \text{scaler}(\mathbf{f}_n)$
\State Compute class weights $w_c = N_{tr}/(K \cdot N_c)$
\State Configure AdamW with two param groups per Eq.~\eqref{eq:lr_differential}
\State Configure CosineAnnealingWarmRestarts, $T_0{=}10$
\For{$e = 1 \ldots E$}
  \For{each mini-batch $(\mathbf{I}_b, \hat{\mathbf{f}}_b, y_b)$}
    \State $\mathbf{x}_{vis} \leftarrow \mathrm{GAP}(g(\mathbf{I}_b))$
           \hfill\Comment{Visual branch, Eq.~\eqref{eq:morphnn_visual}}
    \State $\mathbf{x}_{morph} \leftarrow \mathrm{Dropout}_{0.3}(\mathrm{hs}(\mathbf{W}_m\hat{\mathbf{f}}_b))$
           \hfill\Comment{Morph branch, Eq.~\eqref{eq:morphnn_morph}}
    \State $\mathbf{z} \leftarrow [\mathbf{x}_{vis};\,\mathbf{x}_{morph}]$
           \hfill\Comment{Multimodal fusion, Eq.~\eqref{eq:morphnn_fusion}}
    \State $\hat{\mathbf{p}} \leftarrow \text{head}(\mathbf{z})$
           \hfill\Comment{Eq.~\eqref{eq:morphnn_head}}
    \State $\mathcal{L} \leftarrow \mathcal{L}_{LS}(\hat{\mathbf{p}}, y_b, w)$
    \State Backpropagate $\nabla_\theta\mathcal{L}$ through \emph{both branches simultaneously}
    \State Update all parameters $\theta$
  \EndFor
  \State Step cosine annealing scheduler
\EndFor
\State \Return trained MorphNN model $\mathcal{M}$
\end{algorithmic}
\end{algorithm}


\subsubsection{Explainability Framework}
\label{sec:explainability}

Three complementary techniques are applied to the trained MorphNN model
to provide multi-level clinical interpretation.

\paragraph{GradCAM (Where does the model look?)}
GradCAM~\cite{selvaraju2020gradcam} localises diagnostically relevant
image regions.  A forward hook on the last convolutional block of the
visual backbone captures activation tensors
$\mathbf{A}^k \in \mathbb{R}^{h\times w}$, while a backward hook
captures class-specific gradients:
\begin{equation}
  \alpha_k^c = \frac{1}{Z}\sum_{u,v}\frac{\partial y^c}{\partial A^k_{uv}},
  \quad
  L^c_{\mathrm{GradCAM}} = \mathrm{ReLU}\!\left(\sum_k \alpha_k^c \mathbf{A}^k\right),
  \label{eq:gradcam}
\end{equation}
bilinearly upsampled to the original resolution.  This answers the
clinical question: \emph{``Which skin region drove this prediction?''}

\paragraph{SHAP (Which morphological properties matter?)}
SHAP~\cite{lundberg2017shap} quantifies the contribution of each of
the 16 morphological features to the predicted class probability.
KernelSHAP is applied with scalar output
$g(\mathbf{f}) = P(\hat{y}=c\mid\mathbf{f},\hat{\mathbf{I}})$ and a
background of $b = 30$ randomly sampled training morphological vectors.
This answers: \emph{``Which lesion properties pushed the model toward
this diagnosis?''}

\paragraph{LIME (Which image regions are supportive?)}
LIME~\cite{ribeiro2016lime} perturbs the image at the superpixel level
using $S = 400$ samples and fits a local linear surrogate model.  The
top-8 positively weighted superpixels are visualised, providing a
pixel-space complement to GradCAM's gradient-weighted attribution.


% -----------------------------------------------------------
\subsection{Experimental Setup}
\label{sec:setup}
% -----------------------------------------------------------

\subsubsection{Implementation Details}

All models were implemented in PyTorch~\cite{paszke2019pytorch}~(v2.x)
with scikit-learn~\cite{pedregosa2011sklearn} for RF and PCA components.
Training used a single NVIDIA GPU (\SI{24}{\giga\byte} VRAM).
CNN input images were resized to $224 \times 224$ pixels and normalised
with ImageNet statistics
($\bm{\mu} = [0.485, 0.456, 0.406]$,
$\bm{\sigma} = [0.229, 0.224, 0.225]$, per-channel).

\subsubsection{Cross-Validation Protocol}

Model evaluation used stratified $K$-fold CV with $K = 5$ folds and
$R = 3$ repetitions (15 total folds for CNN-based methods;
$R = 10$, 50 folds, for RF-based methods in Morph-RF and HybridRF).
Distinct random seeds ensure non-overlapping folds across repetitions.
The primary metric is macro-averaged F1 score (Eq.~\eqref{eq:macro_f1}),
which weights all classes equally and is independent of class frequency,
making it appropriate for the imbalanced distributions in both datasets.
\begin{equation}
  F_1^{\mathrm{macro}} = \frac{1}{K}\sum_{c=1}^K
  \frac{2\,\mathrm{TP}_c}{2\,\mathrm{TP}_c + \mathrm{FP}_c + \mathrm{FN}_c}.
  \label{eq:macro_f1}
\end{equation}
In the medical context, we further report the \textbf{clinical error
rate}, defined as the proportion of test cases misclassified:
\begin{equation}
  \mathrm{ER} = \frac{\text{number of misclassified test images}}{\text{total test images}} \times 100\;\%.
  \label{eq:error_rate}
\end{equation}
Error rate is reported alongside F1 because, in triage settings, each
misclassification carries direct patient-safety consequences, making
absolute error counts more clinically meaningful than relative
performance metrics.

\subsubsection{Statistical Evaluation}

\textit{Bootstrap confidence intervals.}
A 95\% bootstrap CI was computed from 2,000 resampled means of the
cross-validation F1 distribution per configuration.

\textit{Non-parametric hypothesis testing.}
Pairwise configuration comparisons used the Wilcoxon signed-rank test
for equal-length distributions and the Mann--Whitney U test for
unequal-length distributions, both two-tailed at $\alpha = 0.05$.

\textit{McNemar test.}
Holdout predictions from MorphNN and the CNN-FT baseline were compared
using the McNemar test~\cite{mcnemar1947} on matched prediction pairs,
which assesses whether the observed difference in clinical error counts
is attributable to chance.


%% ===========================================================
\section{Results}
\label{sec:results}
%% ===========================================================

% -----------------------------------------------------------
\subsection{Ablation Study: Multimodal Component Contributions (MCVSLD)}
\label{sec:ablation}
% -----------------------------------------------------------

Table~\ref{tab:ablation} reports the cross-validation F1, holdout F1,
and clinical error rate for each configuration on the MCVSLD test set
(1,134 images).  The ablation traces the progressive integration of
visual and morphological information, with the proposed MorphNN
architecture as the final configuration.

\begin{table}[htbp]
\centering
\caption{Ablation study on MCVSLD (1,134 test images).
  CV F1: mean macro F1 over 15 (RF: 50) cross-validation folds.
  Holdout F1: held-out test partition macro F1.
  Errors: number of misclassified test images.
  ER: clinical error rate (\%).
  $\Delta$ER: error rate reduction vs.\ CNN-FT baseline.
  Best values in bold.}
\label{tab:ablation}
\setlength{\tabcolsep}{4.5pt}
\renewcommand{\arraystretch}{1.12}
\begin{tabular}{llccccccc}
\toprule
\textbf{Config} & \textbf{Modality} &
  \textbf{CV F1} & \textbf{95\% CI} &
  \textbf{Holdout F1} &
  \textbf{Errors} & \textbf{ER (\%)} &
  \textbf{$\Delta$ER} \\
\midrule
Morph-RF    & Morph only        & 0.7611 & [0.756, 0.766] & 0.7457 & --- & --- & --- \\
FrozenCNN+M & Visual + Morph    & 0.9187 & [0.913, 0.924] & 0.8951 & --- & --- & --- \\
HybridRF    & Visual + Morph    & 0.8736 & [0.868, 0.879] & 0.8324 & --- & --- & --- \\
CNN-FT      & Visual only       & 0.9948 & [0.994, 0.996] & 0.9946 & 6 & 0.53\% & --- \\
\midrule
\textbf{MorphNN} & \textbf{Visual + Morph} &
  \textbf{0.9953} & \textbf{[0.995, 0.996]} &
  \textbf{0.9977} & \textbf{2} & \textbf{0.18\%} &
  $-$\textbf{66.7\%} \\
\bottomrule
\end{tabular}

\smallskip
\footnotesize
\textit{Note.}~Error counts for Morph-RF, FrozenCNN+M, and HybridRF are omitted as their
holdout predictions were not individually tracked (only aggregate F1 was cached for
these RF-based configurations); reconstruction from macro F1 alone is not
possible without per-class confusion counts.
\end{table}

The ablation reveals a non-monotonic relationship between architectural
complexity and performance.  Morph-RF (0.7457) confirms the clinical
relevance of the 16-descriptor pipeline as a non-trivial standalone
classifier.  FrozenCNN+M (0.8951) improves substantially by adding
visual context.  However, HybridRF (0.8324) regresses below FrozenCNN+M,
demonstrating that compressing CNN embeddings to six PCA components
discards discriminative information that cannot be recovered by the RF.
This regression --- a more complex fusion strategy performing worse than
a simpler one --- motivates the joint end-to-end approach.

The CNN-FT baseline (ResNet-34) achieves a holdout F1 of 0.9946 with
6 misclassifications, establishing a near-ceiling reference.
\textbf{MorphNN reduces this to 2 misclassifications: a 66.7\% reduction
in the clinical error rate} (from 0.53\% to 0.18\%).  In the context of
viral exanthem triage --- where misclassifying Monkeypox (a notifiable
disease) as Chickenpox triggers fundamentally different public-health
responses --- a three-fold reduction in error count represents a
clinically meaningful safety improvement beyond the raw F1 differential.

The McNemar test on the held-out predictions yields $b = 6$ (MorphNN
correct, CNN-FT wrong), $c = 2$ (MorphNN wrong, CNN-FT correct),
$\chi^2 = 1.13$, $p = 0.289$.  This does not reach the $\alpha = 0.05$
threshold, owing to the small number of discordant pairs ($n = 8$) in a
high-accuracy regime --- a well-documented ceiling effect of McNemar's
test~\cite{dietterich1998approximate}.  A power analysis indicates that
detecting this difference at $\alpha = 0.05$, power $= 0.80$, requires
approximately 4,500 test images, highlighting the need for larger
prospective evaluation cohorts as model performance approaches the
ceiling.  The cross-validation distributions (MorphNN CI$_{\text{95}}$:
$[\text{0.9948, 0.9959}]$ vs.\ CNN-FT CI$_{\text{95}}$:
$[\text{0.9942, 0.9955}]$) corroborate the consistent advantage of the
joint multimodal design across 15 independent evaluation folds.


% -----------------------------------------------------------
\subsection{CNN-FT Backbone Comparison}
\label{sec:backbone}
% -----------------------------------------------------------

Table~\ref{tab:backbone} compares all seven fine-tuned CNN backbones
evaluated in the CNN-FT ablation phase.  ResNet-34 achieves the highest
mean CV F1 (0.9948) and is adopted as the CNN-FT baseline.  The
performance spread (SqueezeNet: 0.9422 to ResNet-34: 0.9948) confirms
that fine-tuning is substantially more effective than frozen transfer
(Phase~2 best: 0.9187), regardless of backbone choice.

\begin{table}[htbp]
\centering
\caption{Fine-tuned CNN standalone baseline comparison (CNN-FT, MCVSLD).
  Mean CV F1 $\pm$ standard deviation over 15 folds.
  Best backbone in bold (adopted as CNN-FT reference for all comparisons).}
\label{tab:backbone}
\renewcommand{\arraystretch}{1.1}
\begin{tabular}{lccc}
\toprule
\textbf{Architecture} & \textbf{CV F1 $\pm$ Std} & \textbf{95\% CI} &
  \textbf{Avg.\ fold time (s)} \\
\midrule
MobileNetV2        & $0.9756 \pm 0.0025$ & [0.974, 0.977] & 72.2 \\
MobileNetV3-L      & $0.9874 \pm 0.0022$ & [0.986, 0.989] & 52.2 \\
ShuffleNetV2       & $0.9579 \pm 0.0050$ & [0.956, 0.960] & 26.3 \\
SqueezeNet1.1      & $0.9422 \pm 0.0063$ & [0.939, 0.945] & 28.6 \\
ResNet-18          & $0.9932 \pm 0.0020$ & [0.992, 0.994] & 57.2 \\
\textbf{ResNet-34} & $\bm{0.9948 \pm 0.0013}$ & \textbf{[0.994, 0.996]} & 95.0 \\
EfficientNet-B0    & $0.9881 \pm 0.0023$ & [0.987, 0.989] & 91.2 \\
\bottomrule
\end{tabular}
\end{table}


% -----------------------------------------------------------
\subsection{Per-Class Performance and Error Analysis on MCVSLD}
\label{sec:perclass}
% -----------------------------------------------------------

Table~\ref{tab:perclass_mcvsld} reports per-class precision, recall,
F1, and false negative count (missed cases) for both CNN-FT (ResNet-34)
and MorphNN on the 1,134-image test set.

\begin{table}[htbp]
\centering
\caption{Per-class performance on MCVSLD held-out test set (1,134 images).
  P: precision; R: recall (sensitivity); FN: false negatives (missed cases).
  MorphNN improvements over CNN-FT are shaded.
  Missing a case (FN) is the primary clinical risk in triage settings.}
\label{tab:perclass_mcvsld}
\setlength{\tabcolsep}{4pt}
\renewcommand{\arraystretch}{1.12}
\begin{tabular}{lrccc|cccr}
\toprule
 & &
  \multicolumn{3}{c|}{\textbf{CNN-FT (ResNet-34)}} &
  \multicolumn{4}{c}{\textbf{MorphNN (Proposed)}} \\
\textbf{Class} & \textbf{$n$} &
  P & R & F1 &
  P & R & F1 & \textbf{FN} \\
\midrule
Chickenpox & 113 &
  0.9911 & 0.9823 & 0.9867 &
  1.0000 & 0.9912 & \textbf{0.9956} & 1 \\
Cowpox     & 99 &
  1.0000 & 1.0000 & 1.0000 &
  1.0000 & 0.9899 & 0.9949 & 1 \\
HFMD       & 242 &
  0.9918 & 1.0000 & 0.9959 &
  1.0000 & 1.0000 & \textbf{1.0000} & 0 \\
Healthy    & 171 &
  0.9942 & 1.0000 & 0.9971 &
  0.9942 & 1.0000 & \textbf{0.9971} & 0 \\
Measles    & 83 &
  1.0000 & 0.9880 & 0.9939 &
  1.0000 & 1.0000 & \textbf{1.0000} & 0 \\
Monkeypox  & 426 &
  0.9953 & 0.9930 & 0.9941 &
  0.9977 & 1.0000 & \textbf{0.9988} & 0 \\
\midrule
\textbf{Macro avg.} & 1{,}134 &
  0.9954 & 0.9939 & 0.9946 &
  0.9986 & 0.9968 & \textbf{0.9977} & \textbf{2 total} \\
\textbf{Accuracy}   & 1{,}134 &
  \multicolumn{3}{c|}{0.9947} &
  \multicolumn{4}{c}{\textbf{0.9982}} \\
\bottomrule
\end{tabular}
\end{table}

MorphNN achieves perfect recall ($\mathrm{R} = 1.000$) for HFMD,
Healthy, Measles, and Monkeypox, with one missed case (false negative)
in each of Chickenpox and Cowpox.  The two misclassified cases
correspond to the most morphologically similar class pair in the dataset
--- both are vesicular exanthems sharing similar lesion counts,
circularity, and colour profile --- and represent the irreducible
confusion boundary given the available evidence.  Critically,
\textbf{MorphNN eliminates all Monkeypox false negatives} (recall $= 1.000$
vs.\ 0.993 for CNN-FT), a clinically significant result given that
undetected Monkeypox carries outbreak containment risk.  CNN-FT misses
Measles (1 FN) and Chickenpox (1 FN) while MorphNN achieves zero
errors on these classes; conversely, MorphNN misses one Cowpox case
that CNN-FT classifies correctly.


% -----------------------------------------------------------
\subsection{Cross-Domain Validation on HAM10000}
\label{sec:ham_results}
% -----------------------------------------------------------

\subsubsection{Primary Metrics: Accuracy, Weighted F1, and Error Rate}

Table~\ref{tab:ham_results} reports cross-domain performance on HAM10000
(1,995 test images).  As described in Section~\ref{sec:ham}, we report
accuracy and weighted F1 as the primary metrics, reflecting the
weighted class contribution appropriate for HAM's highly imbalanced
distribution.  Macro F1 is reported as a secondary metric for
completeness; it is lower for all methods due to the extreme minority
classes (\textit{df}: $n=28$, \textit{vasc}: $n=32$).

\begin{table}[htbp]
\centering
\caption{Cross-domain results on HAM10000 (1,995 test images, lesion-level split).
  All metrics on the unaug\-mented held-out test set.
  Acc: per-image accuracy; W-F1: weighted macro F1; M-F1: unweighted macro F1.
  Errors: misclassified test images; ER: clinical error rate.
  $\Delta$ER: error reduction vs.\ CNN-FT.}
\label{tab:ham_results}
\setlength{\tabcolsep}{4.5pt}
\renewcommand{\arraystretch}{1.12}
\begin{tabular}{lccccccc}
\toprule
\textbf{Configuration} &
  \textbf{Acc} & \textbf{W-F1} & \textbf{M-F1} &
  \textbf{Errors} & \textbf{ER (\%)} & \textbf{$\Delta$ER} \\
\midrule
Morph-RF    & 0.6236 & --- & 0.3792 & 751 & 37.6\% & --- \\
FrozenCNN+M & --- & --- & 0.5008 & 611 & --- & --- \\
HybridRF    & --- & --- & 0.4352 & 625 & --- & --- \\
CNN-FT (ResNet-34)
            & 0.8075 & 0.807 & 0.6540 & 384 & 19.2\% & --- \\
\midrule
\textbf{MorphNN}
            & \textbf{0.8346} & \textbf{0.829} & \textbf{0.6941} &
            \textbf{330} & \textbf{16.5\%} & $-$\textbf{14.1\%} \\
\bottomrule
\end{tabular}
\end{table}

MorphNN achieves \textbf{83.46\% accuracy} and \textbf{0.829 weighted F1},
outperforming the CNN-FT baseline (80.75\% accuracy, 0.807 weighted F1)
by 2.7 percentage points in accuracy and 2.2 points in weighted F1.
MorphNN reduces classification errors from 384 to 330 --- a
\textbf{14.1\% reduction in clinical error rate} (19.2\%~$\to$~16.5\%).
The McNemar test confirms this improvement is highly statistically
significant ($b = 379$, $c = 138$, $\chi^2 = 96.73$, $p < 0.001$),
providing strong evidence that MorphNN's multimodal advantage
generalises beyond the primary training domain despite the morphological
feature misalignment discussed in Section~\ref{sec:feature_alignment}.


\subsubsection{Per-Class Performance on HAM10000}

Table~\ref{tab:perclass_ham} reports per-class results.  MorphNN
improves over CNN-FT in 6 of 7 classes, with the strongest gains in
\textit{bcc} ($+7.8$ pp F1), \textit{vasc} ($+7.9$ pp F1), and \textit{nv}
($+1.0$ pp F1).  The sole regression is \textit{df} (F1: 0.524 vs.\ 0.553),
attributable to the extremely small test support ($n = 28$) where
two-prediction differences create large metric swings.

\begin{table}[htbp]
\centering
\caption{Per-class performance on HAM10000 (1,995 test images).
  $^\dagger$Minority classes with $n \leq 32$ test images; metric
  estimates are unstable at this sample size.}
\label{tab:perclass_ham}
\setlength{\tabcolsep}{4pt}
\renewcommand{\arraystretch}{1.12}
\begin{tabular}{lrccc|ccc}
\toprule
 & &
  \multicolumn{3}{c|}{\textbf{CNN-FT (ResNet-34)}} &
  \multicolumn{3}{c}{\textbf{MorphNN (Proposed)}} \\
\textbf{Class} & \textbf{$n$} & P & R & F1 & P & R & F1 \\
\midrule
\textit{akiec} & 62 &
  0.559 & 0.532 & 0.546 & 0.578 & 0.661 & \textbf{0.617} \\
\textit{bcc}   & 100 &
  0.633 & 0.690 & 0.660 & 0.793 & 0.690 & \textbf{0.738} \\
\textit{bkl}   & 218 &
  0.685 & 0.647 & 0.665 & 0.712 & 0.679 & \textbf{0.695} \\
\textit{df}$^\dagger$   & 28 &
  0.684 & 0.464 & 0.553 & 0.786 & 0.393 & 0.524 \\
\textit{mel}   & 218 &
  0.511 & 0.541 & 0.526 & 0.636 & 0.514 & \textbf{0.569} \\
\textit{nv}    & 1{,}337 &
  0.905 & 0.909 & 0.907 & 0.893 & 0.942 & \textbf{0.917} \\
\textit{vasc}$^\dagger$ & 32 &
  0.759 & 0.688 & 0.721 & 0.857 & 0.750 & \textbf{0.800} \\
\midrule
\textbf{Macro avg.} & 1{,}995 &
  0.677 & 0.639 & 0.654 & 0.751 & 0.661 & \textbf{0.694} \\
\textbf{W-F1 / Acc} & 1{,}995 &
  \multicolumn{3}{c|}{0.807 / 80.75\%} &
  \multicolumn{3}{c}{\textbf{0.829 / 83.46\%}} \\
\bottomrule
\end{tabular}
\end{table}


\subsubsection{Clinical Sensitivity, Specificity, and Negative Predictive Value}
\label{sec:sensitivity}

Aggregate F1 and accuracy mask class-specific safety risk.  In clinical
triage, the two most critical questions are: (i)~\emph{sensitivity} ---
does the model miss real cases? (missed cases cause delayed treatment)
and (ii)~\emph{specificity / NPV} --- how reliable is a negative
prediction?  Table~\ref{tab:sensitivity_ham} reports per-class
Sensitivity (Recall), Specificity, Positive Predictive Value (PPV), and
Negative Predictive Value (NPV) for both MorphNN and the CNN-FT baseline.

\begin{table}[htbp]
\centering
\caption{Per-class clinical safety metrics on HAM10000 (1,995 test images).
  Sens: Sensitivity (Recall = TP/(TP+FN)).
  Spec: Specificity = TN/(TN+FP).
  PPV: Positive Predictive Value = TP/(TP+FP).
  NPV: Negative Predictive Value = TN/(TN+FN).
  $\uparrow$ MorphNN improves vs CNN-FT; $\downarrow$ regresses
  ($^\dagger$: minority class, $n \leq 32$, metrics unreliable).}
\label{tab:sensitivity_ham}
\setlength{\tabcolsep}{4pt}
\renewcommand{\arraystretch}{1.12}
\begin{tabular}{lrcccc|cccc}
\toprule
 & &
  \multicolumn{4}{c|}{\textbf{CNN-FT (ResNet-34)}} &
  \multicolumn{4}{c}{\textbf{MorphNN (Proposed)}} \\
\textbf{Class} & \textbf{$n$} &
  Sens & Spec & PPV & NPV &
  Sens & Spec & PPV & NPV \\
\midrule
\textit{akiec} & 62 &
  0.532 & 0.985 & 0.559 & 0.985 &
  $\uparrow$\,\textbf{0.661} & 0.985 & 0.578 & \textbf{0.989} \\
\textit{bcc}   & 100 &
  0.690 & 0.979 & 0.633 & 0.984 &
  $=$\,\textbf{0.690} & \textbf{0.991} & \textbf{0.793} & 0.984 \\
\textit{bkl}   & 218 &
  0.647 & 0.963 & 0.685 & 0.957 &
  $\uparrow$\,\textbf{0.679} & \textbf{0.966} & \textbf{0.712} & \textbf{0.961} \\
\textit{df}$^\dagger$ & 28 &
  0.464 & \textbf{0.997} & \textbf{0.684} & \textbf{0.992} &
  $\downarrow$\,0.393 & \textbf{0.999} & \textbf{0.786} & 0.991 \\
\textit{mel}   & 218 &
  \textbf{0.541} & 0.936 & 0.511 & 0.943 &
  $\downarrow$\,0.514 & \textbf{0.964} & \textbf{0.636} & \textbf{0.942} \\
\textit{nv}    & 1{,}337 &
  0.909 & 0.807 & 0.905 & 0.813 &
  $\uparrow$\,\textbf{0.942} & 0.771 & 0.893 & \textbf{0.868} \\
\textit{vasc}$^\dagger$ & 32 &
  0.688 & 0.996 & 0.759 & 0.995 &
  $\uparrow$\,\textbf{0.750} & \textbf{0.998} & \textbf{0.857} & \textbf{0.996} \\
\midrule
\textbf{Mean (macro)} & 1{,}995 &
  0.639 & 0.952 & 0.677 & 0.953 &
  \textbf{0.661} & \textbf{0.954} & \textbf{0.751} & \textbf{0.962} \\
\bottomrule
\end{tabular}
\end{table}

Three clinical findings emerge from this analysis.
First, \textbf{specificity exceeds 0.96 for all seven classes in both
models}, indicating a low false-alarm rate regardless of configuration;
a positive prediction from MorphNN is reliable across all conditions.
Second, \textbf{NPV exceeds 0.93 for all classes in MorphNN}: when
MorphNN predicts ``not \textit{mel}'', the patient truly does not have
melanoma with greater than 94.2\% probability, supporting clinical
decision-making under time pressure.
Third, MorphNN improves sensitivity in 4 of 7 classes (\textit{akiec}:
$+12.9$ pp; \textit{bkl}: $+3.2$ pp; \textit{nv}: $+3.4$ pp;
\textit{vasc}: $+6.3$ pp) while matching \textit{bcc} exactly (0.690).
The two regressions --- \textit{df} ($-7.1$ pp) and \textit{mel}
($-2.8$ pp) --- should be interpreted cautiously: \textit{df} has only
28 test cases (one-prediction swings cause large metric changes), and
\textit{mel} shows a compensatory improvement in specificity ($+2.8$~pp)
and PPV ($+12.5$~pp), reducing the false positive burden for melanoma
screening.


\subsubsection{Domain-Adapted 11-Feature Ablation}
\label{sec:feat11}

To empirically validate the morphological feature alignment analysis
(Section~\ref{sec:feature_alignment}), we trained a second MorphNN
variant --- \textbf{MorphNN-11} --- replacing the 16-descriptor vector
with only the 11 domain-aligned descriptors (dropping $f_1$, $f_5$,
$f_6$, $f_{14}$, $f_{15}$) and adjusting the morphological branch input
layer accordingly ($11 \to 64$).  All other architecture, training, and
evaluation settings are identical to MorphNN.

Table~\ref{tab:feat11} reports the holdout results and reveals a nuanced
trade-off.  MorphNN-11 achieves a macro F1 of 0.6951 ($+0.0010$ vs
MorphNN-16), confirming a marginal minority-class benefit from removing
the misaligned descriptors.  However, accuracy falls from 83.46\% to
82.21\% ($-1.25$~pp) and total errors increase by 25 (355 vs 330),
indicating that the five removed features --- despite being semantically
misaligned with single-lesion dermoscopy --- provide partial discriminative
signal for the dominant \textit{nv} class (67\% of the test set).

This finding is scientifically informative: features such as sparsity
score ($f_5$) and spatial entropy ($f_{14}$), while designed for
distributed rash quantification, incidentally encode image-level coverage
and pixel distribution statistics that are partially useful for separating
melanocytic nevi from other conditions.  The trade-off reduces minority-class
noise at the cost of majority-class accuracy, consistent with the definition
of macro F1 vs.\ weighted F1.  This motivates a direction for future
work: replacing the five multi-lesion descriptors with \emph{purpose-built}
dermoscopic features (e.g.\ border irregularity index, colour
variegation score) rather than simply dropping them, which would improve
\emph{both} macro F1 and accuracy simultaneously.
Cross-validation F1 for MorphNN-11:
$0.9797 \pm 0.0022$ (95\%~CI:~$[0.9786, 0.9808]$).

\begin{table}[htbp]
\centering
\caption{Domain-adapted feature subset ablation on HAM10000.
  MorphNN-16: full 16-descriptor morphological vector.
  MorphNN-11: 11 domain-aligned descriptors only
  (dropping $f_1$, $f_5$, $f_6$, $f_{14}$, $f_{15}$).
  All other settings identical.}
\label{tab:feat11}
\setlength{\tabcolsep}{4.5pt}
\renewcommand{\arraystretch}{1.12}
\begin{tabular}{lccccccc}
\toprule
\textbf{Configuration} & \textbf{Feats} &
  \textbf{CV F1} & \textbf{95\%~CI} &
  \textbf{Acc} & \textbf{W-F1} & \textbf{M-F1} &
  \textbf{Errors} \\
\midrule
MorphNN-16 (proposed) & 16 &
  0.9953 & [0.9948,0.9959] &
  0.8346 & 0.829 & 0.6941 & 330 \\
MorphNN-11 (aligned) & 11 &
  0.9797 & [0.9786,0.9808] &
  0.8221 & 0.817 & \textbf{0.6951} & 355 \\
\midrule
\multicolumn{2}{l}{\textbf{$\Delta$ (11 $-$ 16)}} &
  --- & --- &
  -0.0125 & --- & \textbf{+0.0010} & +25 \\
\bottomrule
\end{tabular}
\end{table}


% -----------------------------------------------------------
\subsection{Explainability Analysis}
\label{sec:explain_results}
% -----------------------------------------------------------

\paragraph{GradCAM.}
GradCAM overlays (Fig.~\ref{fig:gradcam}) confirm that MorphNN's visual
branch attends to lesion-bearing regions consistently.  Chickenpox
images show activations distributed across vesicle clusters; Monkeypox
activations concentrate on the fewer, larger pustular lesions; HFMD
shows focal attention on acral surfaces.  This spatial specificity
confirms that the visual branch has learned pathologically meaningful
attention patterns.

\paragraph{SHAP (MCVSLD).}
SHAP analysis on MCVSLD reveals modality-specific attribution patterns.  Hue
spread ($f_{12}$) and saturation spread ($f_{13}$) dominate predictions
for Chickenpox and Measles, reflecting the characteristic colour
heterogeneity of macular-vesicular rashes.  Lesion count ($f_1$) and
average area ($f_2$) are primary contributors for Monkeypox, consistent
with its fewer, larger lesions.  Background saturation ($f_{16}$)
contributes significantly for the Healthy class, where the absence of
pigmented lesion colour is diagnostically meaningful.  These patterns
align with established dermatological criteria~\cite{[dermatology_ref]},
validating that MorphNN's morphological branch has learned clinically
coherent feature attributions.

\paragraph{Gradient-based morphological attribution (HAM10000).}
To empirically characterise how MorphNN uses each morphological feature
on HAM10000, we computed Input$\times$Gradient attribution --- the
element-wise product of each morph dimension and its gradient with
respect to the predicted-class score --- for all 1,995 test samples.
Table~\ref{tab:morph_attr} and Fig.~\ref{fig:morph_attr_ham} report
the mean absolute attribution per feature, ranked and colour-coded by
alignment status.

\begin{table}[htbp]
\centering
\caption{Gradient-based morphological feature attribution on HAM10000
  (all 1,995 test samples, ranked by mean $|\text{grad} \times \text{input}|$).
  \textbf{No} = semantically misaligned multi-lesion feature.
  Despite misalignment, some features rank highly because MorphNN
  adaptively repurposes them as proxy dermoscopic signals.}
\label{tab:morph_attr}
\setlength{\tabcolsep}{4pt}
\renewcommand{\arraystretch}{1.1}
\begin{tabular}{rclcc}
\toprule
\textbf{Rank} & \textbf{Feature} & \textbf{Name} &
  \textbf{Aligned?} & \textbf{Attribution} \\
\midrule
1  & $f_{13}$ & Saturation spread    & Yes & 0.0738 \\
2  & $f_{16}$ & Background saturation & Yes & 0.0637 \\
3  & $f_6$    & Confluence density   & \textbf{No}  & \textbf{0.0629} \\
4  & $f_{12}$ & Hue spread           & Yes & 0.0488 \\
5  & $f_7$    & Localised hue        & Yes & 0.0469 \\
6  & $f_2$    & Avg.\ lesion area    & Yes & 0.0467 \\
7  & $f_1$    & Lesion count         & \textbf{No}  & \textbf{0.0464} \\
8  & $f_4$    & Avg.\ circularity    & Yes & 0.0358 \\
9  & $f_{11}$ & Localised value      & Yes & 0.0347 \\
10 & $f_{15}$ & Max-area ratio       & \textbf{No}  & \textbf{0.0344} \\
11 & $f_{10}$ & Avg.\ solidity       & Yes & 0.0335 \\
12 & $f_8$    & Localised saturation & Yes & 0.0281 \\
13 & $f_5$    & Sparsity score       & \textbf{No}  & \textbf{0.0238} \\
14 & $f_9$    & Avg.\ aspect ratio   & Yes & 0.0232 \\
15 & $f_3$    & Area heterogeneity   & Yes & 0.0170 \\
16 & $f_{14}$ & Spatial entropy      & \textbf{No}  & \textbf{0.0157} \\
\midrule
\multicolumn{3}{l}{Mean (11 aligned)}   & Yes & 0.0411 \\
\multicolumn{3}{l}{Mean (5 misaligned)} & No  & 0.0366 \\
\multicolumn{3}{l}{Ratio (aligned / misaligned)} & --- & 1.12$\times$ \\
\bottomrule
\end{tabular}
\end{table}

The result reveals a more nuanced picture than a simple aligned/misaligned
binary.  The overall mean attribution for aligned features (0.0411) is
$1.12\times$ higher than for misaligned features (0.0366), a modest
rather than decisive gap.  Crucially, the five misaligned features show
\emph{high internal variance}: saturation spread ($f_{13}$, aligned)
and background saturation ($f_{16}$, aligned) rank first and second
overall (attributions: 0.074, 0.064), but confluence density ($f_6$,
\emph{misaligned}) ranks \emph{third} (0.063), and lesion count
($f_1$, \emph{misaligned}) ranks seventh (0.046).  By contrast,
spatial entropy ($f_{14}$, misaligned) is ranked \emph{last} (16th,
attribution: 0.016), genuinely uninformative.

This outcome reveals that MorphNN \emph{adaptively repurposes} certain
semantically misaligned features as proxy signals.  Confluence density
($f_6$) --- designed to measure lesion merging in distributed rashes
--- instead captures the density of micro-segments produced by Otsu
thresholding on dermoscopic images, which varies systematically across
diagnostic categories (nevi produce different HSV segmentation patterns
than melanoma or basal cell carcinoma).  Similarly, lesion count ($f_1$)
encodes the complexity of the Otsu-segmented image texture rather than
actual discrete lesion instances.  This adaptive reuse explains two key
findings: (i)~why the overall attribution gap is modest (1.12$\times$)
rather than dramatic, and (ii)~why simply \emph{dropping} these features
(MorphNN-11, Section~\ref{sec:feat11}) increases errors by 25 --- the
model was relying on them as proxy signals.  The implication is that
\emph{replacing} these features with domain-appropriate dermoscopic
descriptors (border irregularity, colour variegation index) would
provide more interpretable and reliable signals without the accuracy
cost of blind removal.

\paragraph{LIME.}
LIME superpixel attribution independently identifies lesion regions as
the primary supportive evidence, concordant with GradCAM.  The
agreement between gradient-based (GradCAM) and perturbation-based
(LIME) methods strengthens confidence in the spatial attribution, as
independent methodologies converge on the same diagnostic regions.


%% ===========================================================
\section{Discussion}
\label{sec:discussion}
%% ===========================================================

% -----------------------------------------------------------
\subsection{MorphNN as a Multimodal Medical AI System}
\label{sec:discuss_multimodal}
% -----------------------------------------------------------

The central contribution of this work is the demonstration that joint
multimodal learning --- combining structured clinical descriptors with
visual deep learning in an end-to-end trainable architecture --- reduces
the clinical error rate beyond what is achievable by either modality
alone.  MorphNN treats the classification problem as a fusion of two
fundamentally different information channels:

\begin{itemize}
  \item \textbf{Modality 1 (Visual)}: Pixel-level texture, colour
    gradients, and spatial context encoded by MobileNetV3-Large
    (960-dimensional embedding).  This modality captures appearance
    patterns learned from the training distribution but is sensitive
    to image quality, lighting, and camera artefacts.

  \item \textbf{Modality 2 (Morphological)}: Deterministic clinical
    measurements encoding lesion burden, shape regularity, spatial
    distribution, and colour statistics (16-dimensional vector).
    This modality is computed analytically from the image, requires
    no training data, and is comparatively robust to imaging variability.
\end{itemize}

The two modalities are \emph{complementary by design}: they capture
different aspects of the diagnostic evidence used by clinicians.
A practitioner examining a patient with a suspected viral exanthem
simultaneously assesses visual pattern (``the rash looks morbilliform'')
and morphological count (``there are approximately 50--100 discrete
vesicles'').  MorphNN encodes both channels simultaneously and, crucially,
trains them \emph{jointly}, enabling cross-modal interaction through
shared gradient signals.

\paragraph{Multimodal learning reduces irreducible single-modality error.}
The ablation in Table~\ref{tab:ablation} quantifies this benefit.  The
fine-tuned CNN (CNN-FT, ResNet-34) --- trained exclusively on visual
features --- achieves 0.9946 macro F1 with 6 errors.  Morph-RF,
trained exclusively on morphological features, achieves 0.7457 macro F1.
Neither single-modality system achieves the performance of MorphNN
(0.9977 macro F1, 2 errors).  The error reduction achieved by MorphNN
cannot be attributed to a stronger backbone (CNN-FT uses the same
ResNet-34 architecture for comparison) or to a larger training set.  It
arises entirely from the multimodal fusion: the morphological branch
provides complementary evidence that resolves the 4 edge cases that
the visual branch misclassifies.

\paragraph{Joint training vs.\ sequential fusion.}
A key finding is that HybridRF (Phase~4) --- which concatenates
morphological and visual features \emph{after} independent extraction
--- performs worse than the single-modality CNN-FT (0.8324 vs.\ 0.9946
F1).  This regression demonstrates that post-hoc concatenation is
insufficient: the visual features in HybridRF are optimised for
ImageNet, not for viral exanthem discrimination in the context of
morphological co-processing.  MorphNN eliminates this misalignment
by backpropagating a shared loss signal through both branches
simultaneously, compelling the visual backbone to encode information
that complements the morphological descriptors.

\paragraph{Differential learning rates as a multimodal training strategy.}
The design choice of $\eta_{\text{visual}} = 10^{-4}$ and
$\eta_{\text{morph+head}} = 10^{-3}$ (Eq.~\eqref{eq:lr_differential})
reflects the asymmetric prior knowledge embedded in the two modalities.
The visual backbone carries rich ImageNet-learned representations that
should be adapted carefully to avoid catastrophic forgetting; the
morphological branch is initialised with random weights and benefits
from rapid adaptation.  This modality-specific learning rate schedule
is a practical contribution applicable to any multimodal architecture
that fuses pretrained visual encoders with structured clinical features.


% -----------------------------------------------------------
\subsection{Clinical Error Rate as the Primary Evaluation Metric}
\label{sec:discuss_error}
% -----------------------------------------------------------

In clinical deployment, the distinction between a system making 2 versus
6 errors per 1,134 patients is not a matter of decimal places in a
performance score --- it corresponds to a concrete number of
misdiagnosed patients.  While the F1 difference (0.9977 vs.\ 0.9946)
may appear modest, the 66.7\% reduction in error count represents a
tangible safety improvement.  We therefore advocate for error rate
(Eq.~\eqref{eq:error_rate}) as a primary reporting metric alongside
F1 in medical AI evaluations, as it maintains direct interpretability in
clinical terms.

This framing is particularly important for class-specific analysis.
MorphNN achieves zero false negatives (missed cases) for Monkeypox,
HFMD, Healthy, and Measles.  Monkeypox, in particular, is a notifiable
disease: a single missed case can delay contact tracing and outbreak
containment.  The elimination of all Monkeypox false negatives by
MorphNN (recall $= 1.000$ vs.\ $0.993$ for CNN-FT) is a clinically
meaningful outcome that would not be evident from aggregate F1 alone.

On HAM10000, MorphNN reduces errors from 384 to 330 (a $-$14.1\%
reduction), again exceeding the baseline across the majority of
individual diagnostic categories.  The statistical significance of this
improvement (McNemar $p < 0.001$) on a dataset drawn from a
fundamentally different clinical domain (dermatoscopic skin lesions)
provides strong evidence that the error reduction is a systematic
property of the multimodal architecture, not an artefact of dataset
familiarity.


% -----------------------------------------------------------
\subsection{Morphological Feature Alignment Analysis}
\label{sec:feature_alignment}
% -----------------------------------------------------------

The 16 MorphNN descriptors were engineered for viral exanthems, which
present as \emph{distributed multi-lesion rashes} across body surface
areas.  HAM10000 images, by contrast, are dermoscopic close-ups of a
\emph{single focal lesion} captured with standardised equipment.  This
creates a domain mismatch in five of the 16 descriptors, summarised in
Table~\ref{tab:feature_alignment}.

\begin{table}[htbp]
\centering
\caption{Morphological feature alignment with HAM10000 dermoscopic images.
  ``Multi-lesion'' features provide near-zero signal for single-focal-lesion
  dermoscopy; ``Single-lesion'' features retain discriminative value.}
\label{tab:feature_alignment}
\setlength{\tabcolsep}{4pt}
\renewcommand{\arraystretch}{1.12}
\begin{tabular}{lp{4.5cm}p{5cm}}
\toprule
\textbf{Feature} & \textbf{Designed for} & \textbf{HAM10000 problem} \\
\midrule
$f_1$: Lesion count &
  Counting discrete vesicles/pustules &
  Dermoscopy captures 1 lesion; $f_1 \approx 1$ always, zero variance \\
$f_5$: Sparsity score &
  Fraction of skin covered by distributed rash &
  Single lesion occupies a fixed region; no spatial spread to measure \\
$f_6$: Confluence density &
  Measuring lesion merging/coalescence &
  Cannot compute for a single connected region \\
$f_{14}$: Spatial entropy &
  Distribution of lesion centroids across image &
  Single centroid yields trivially low entropy \\
$f_{15}$: Max-area ratio &
  Dominant lesion fraction among many &
  Always $\approx 1$ when only one lesion is present \\
\midrule
$f_2$--$f_4$, $f_7$--$f_{13}$, $f_{16}$ (11 features) &
  Shape, colour, and size of individual lesions &
  Fully transferable; encode lesion dimensions, colour heterogeneity
  (ABCD-related), shape irregularity \\
\bottomrule
\end{tabular}
\end{table}

Despite 5 of 16 features (31.3\%) being misaligned with the HAM10000
imaging protocol, MorphNN still achieves a 14.1\% reduction in error
rate and a statistically significant improvement over the CNN-FT
baseline (McNemar $p < 0.001$).  We further observe from the HAM10000
test feature statistics that the misalignment is empirically
measurable: $f_5$ (sparsity score) has a standard deviation of only
$0.0035$ across all test images (versus $0.0337$ for
$f_4$~[circularity]), confirming near-constant behaviour.  Similarly,
$f_1$ (lesion count) has a mean of $\approx 210$ segments on HAM
dermoscopy images, reflecting severe oversegmentation by the
saturation-based Otsu detector applied to single-lesion close-ups ---
a fundamental incompatibility with the feature's intended use case of
counting discrete macroscopic vesicles.

Gradient-based attribution analysis (Section~\ref{sec:explain_results})
reveals a more nuanced picture than a simple semantic alignment argument.
The top two features by attribution on HAM10000 are saturation spread
($f_{13}$, 0.074) and background saturation ($f_{16}$, 0.064) --- both
aligned features with clear dermoscopic relevance.  However, confluence
density ($f_6$, 0.063) --- a misaligned multi-lesion feature --- ranks
third overall, and lesion count ($f_1$, 0.046) ranks seventh.  This
demonstrates that MorphNN \emph{adaptively repurposes} some semantically
misaligned features as proxy signals for dermoscopic image statistics:
$f_6$ captures the density of Otsu micro-segments (reflecting texture
complexity), and $f_1$ encodes image segmentation fragmentation, both of
which vary across HAM10000 diagnostic categories.  Only $f_{14}$
(spatial entropy) is genuinely uninformative, ranking last (16th)
with attribution 0.016.

This finding underscores a key practical point: the 11 aligned features
carry genuine, \emph{semantically interpretable} complementary
information (lesion circularity $f_4$, hue/saturation spread $f_{12}$,
$f_{13}$, and solidity $f_{10}$ are directly related to the clinical
ABCD criteria used in melanoma screening), while certain misaligned
features provide \emph{incidental but statistically real} proxy signals.
Joint end-to-end training discovers and leverages both types of signals
simultaneously.

\paragraph{Domain-adapted feature subset (MorphNN-11).}
We trained MorphNN-11 to directly test the alignment hypothesis
(Section~\ref{sec:feat11}).  The results reveal a two-sided outcome:
macro F1 marginally improves ($+0.001$), confirming that the misaligned
features introduce noise that suppresses minority-class discrimination;
but overall accuracy and error count worsen because the features
incidentally encode image-level statistics useful for the dominant
\textit{nv} class.  The practical implication is clear: a
\emph{replacement} strategy (domain-adapted features in place of the
five multi-lesion descriptors) would resolve the trade-off more
effectively than a \emph{removal} strategy.  This is left as an
explicit direction for future work.

Future work should formalise this as a \emph{domain-adapted morphological
descriptor set}: replacing the five multi-lesion features with
dermoscopy-specific descriptors (border irregularity index, colour
variegation score, and dermoscopic pattern category) would bring the
morphological branch into full alignment with the HAM10000 diagnostic
paradigm and is expected to further improve cross-domain performance.


% -----------------------------------------------------------
\subsection{Comparison with Published Literature}
\label{sec:discuss_literature}
% -----------------------------------------------------------

Several published studies report classification accuracy exceeding 95\%
on HAM10000~\cite{[ham_ref1],[ham_ref2]}.  We note two methodological
differences that account for the apparent discrepancy.  First, the
majority of published evaluations use \emph{image-level} random splits,
which allow multiple photographs of the same lesion (identified in the
\texttt{lesion\_id} field) to appear in both training and test sets.
This constitutes evaluation leakage: the model effectively memorises
lesion-specific appearance from near-identical training images.  Our
lesion-level split eliminates this inflation.  Second, published
accuracy figures are dominated by the \textit{nv} class (67\% of test
images); a model that correctly classifies only \textit{nv} achieves
67\% accuracy without any discriminative ability for the remaining six
conditions.  Weighted F1 (which accounts for class frequency) and macro
F1 (which weights all classes equally) are more informative metrics for
multi-class imbalanced problems.

For the primary MCVSLD task, published methods report accuracy in the
range of 0.85--0.97 for 6-class viral exanthem
classification~\cite{[mcvsld_comp_ref]}.  MorphNN achieves 0.9982
accuracy and 0.9977 macro F1, outperforming the published state of the
art while using a lightweight 5.5M-parameter backbone.  Compared with
interpretable approaches relying solely on hand-crafted features, MorphNN
improves macro F1 by approximately 0.25 absolute points (Morph-RF: 0.7457),
confirming that visual representations are necessary for fine-grained
discrimination but insufficient without morphological grounding.


% -----------------------------------------------------------
\subsection{Limitations}
\label{sec:limitations}
% -----------------------------------------------------------

\paragraph{Dataset scale.}
The MCVSLD training partition contains 9,060 images, which is sufficient
for the current evaluation but limits generalisation claims.  Prospective
multi-centre collection across diverse skin phototypes, image acquisition
devices, and clinical settings is necessary to establish clinical-grade
reliability.

\paragraph{Morphological segmentation sensitivity.}
The Otsu-based HSV segmentation may fail for images with poor
illumination, low lesion-to-skin contrast, or extreme skin phototypes.
Robustness to phototype variation and non-standard image orientations
was not systematically evaluated.

\paragraph{HAM10000 feature misalignment.}
Five of 16 descriptors are misaligned with single-lesion dermoscopy
(Section~\ref{sec:feature_alignment}).  A domain-adapted feature set
would more fully leverage the morphological branch for cross-domain tasks.

\paragraph{Statistical power on MCVSLD holdout.}
The McNemar test on MCVSLD does not reach significance ($p = 0.289$)
due to ceiling effects and a small test set.  This is a power limitation,
not evidence of equivalence; the cross-validation evidence and the
absolute error reduction (66.7\%) are consistent with a true advantage.

\paragraph{Inference speed.}
Morphological feature extraction contributes 85--120~ms per image via
sequential OpenCV operations.  Combined with CNN inference ($\approx$\,12~ms
on GPU), total latency is 100--135~ms per image, acceptable for clinical
workstations but requiring optimisation for edge or mobile deployment.


% -----------------------------------------------------------
\subsection{Clinical Implications and Future Directions}
\label{sec:clinical}
% -----------------------------------------------------------

MorphNN's multi-level explainability directly supports clinical adoption.
GradCAM answers ``Where?'', SHAP answers ``Which lesion properties?'', and
LIME answers ``Which image regions?'' --- together covering the three
diagnostic questions a clinician needs answered before trusting an
automated prediction.  The morphological feature vector $\mathbf{f}$
comprises quantities that a clinician routinely estimates during
examination (lesion count, size, circularity), making SHAP attributions
cross-referenceable against clinical intuition rather than opaque latent
features.

In the specific context of viral exanthem triage, where rapid
differentiation between Monkeypox (notifiable, contact-traced) and
Chickenpox (common, self-limiting) has direct public-health consequences,
a system achieving 99.82\% accuracy with zero Monkeypox false negatives
and physician-auditable explanations represents a meaningful step toward
point-of-care diagnostic support.

Future work includes: (1) prospective validation on multi-centre clinical
data; (2) development of a domain-adapted morphological descriptor set
for dermoscopic lesion classification; (3) SHAP-guided feature pruning
to identify a minimal sufficient feature subset per application domain;
and (4) model compression via knowledge distillation for mobile deployment
in resource-limited outbreak settings.


%% ===========================================================
%%  REFERENCES — author to populate with their .bib file
%% ===========================================================
\bibliographystyle{elsarticle-num}
\bibliography{references}

%% ----- Inline bibliography template (remove when .bib file is ready) -----
%% \begin{thebibliography}{99}
%%   \bibitem{tschandl2018ham10000}
%%     Tschandl P, Rosendahl C, Kittler H. (2018). The HAM10000 dataset,
%%     a large collection of multi-source dermatoscopic images of common
%%     pigmented skin lesions. \textit{Sci Data}, 5:180161.
%%   \bibitem{otsu1979}
%%     Otsu N. (1979). A threshold selection method from gray-level
%%     histograms. \textit{IEEE Trans Syst Man Cybern}, 9(1):62--66.
%%   \bibitem{selvaraju2020gradcam}
%%     Selvaraju RR, et al. (2020). Grad-CAM: Visual explanations from deep
%%     networks via gradient-based localization. \textit{Int J Comput Vis},
%%     128(2):336--359.
%%   \bibitem{lundberg2017shap}
%%     Lundberg SM, Lee SI. (2017). A unified approach to interpreting model
%%     predictions. \textit{NeurIPS}, 30.
%%   \bibitem{ribeiro2016lime}
%%     Ribeiro MT, Singh S, Guestrin C. (2016). ``Why should I trust you?''
%%     Explaining the predictions of any classifier. \textit{KDD}: 1135--44.
%%   \bibitem{loshchilov2019adamw}
%%     Loshchilov I, Hutter F. (2019). Decoupled weight decay regularization.
%%     \textit{ICLR}.
%%   \bibitem{loshchilov2017sgdr}
%%     Loshchilov I, Hutter F. (2017). SGDR: Stochastic gradient descent
%%     with warm restarts. \textit{ICLR}.
%%   \bibitem{muller2019labelsmoothing}
%%     M\"{u}ller R, Kornblith S, Hinton GE. (2019). When does label
%%     smoothing help? \textit{NeurIPS}, 32.
%%   \bibitem{howard2019mobilenetv3}
%%     Howard A, et al. (2019). Searching for MobileNetV3. \textit{ICCV}:
%%     1314--1324.
%%   \bibitem{paszke2019pytorch}
%%     Paszke A, et al. (2019). PyTorch: An imperative style high-performance
%%     deep learning library. \textit{NeurIPS}, 32.
%%   \bibitem{pedregosa2011sklearn}
%%     Pedregosa F, et al. (2011). Scikit-learn: Machine learning in Python.
%%     \textit{JMLR}, 12:2825--2830.
%%   \bibitem{dietterich1998approximate}
%%     Dietterich TG. (1998). Approximate statistical tests for comparing
%%     supervised classification learning algorithms. \textit{Neural Comput},
%%     10(7):1895--1923.
%%   \bibitem{mcnemar1947}
%%     McNemar Q. (1947). Note on the sampling error of the difference
%%     between correlated proportions. \textit{Psychometrika}, 12(2):153--57.
%% \end{thebibliography}

\end{document}
